\documentclass[aps,prl,twocolumn,superscriptaddress,floatfix]{revtex4-2}
\usepackage{amsmath,amssymb,graphicx,booktabs}
\begin{document}

\title{From-Scratch Replication: Picosecond-Scale Octupole Relaxation in\\
Chiral Antiferromagnets (Konakanchi \textit{et al.}, arXiv:2501.18978)}
\author{Automated Replication Agent}
\affiliation{TEXTURES-100 replication campaign}
\date{\today}

\begin{abstract}
We independently reproduce the central claim of Konakanchi \textit{et al.}
(2025): that the octupole order parameter of a nanoscale chiral
antiferromagnet (Mn$_3$Sn-type cluster octupole) relaxes on
\emph{picosecond}-to-tens-of-picosecond timescales in the low-barrier
regime---orders of magnitude faster than dipolar ferromagnets. Building the
effective low-energy octupole free energy from scratch, we (i) evaluate the
low-barrier precessional-dephasing autocorrelation both analytically and by
Boltzmann sampling, (ii) evaluate the high-barrier Langer/IHD escape time, and
(iii) confirm both with a direct stochastic (Langevin) integration of the
reduced $(m_z,\phi_{\mathrm{oct}})$ dynamics. Our minimum low-barrier
relaxation time is $\tau\simeq 6$--$13$~ps, matching the paper's
``$\sim$10~ps for sub-$kT$ barriers.'' \textbf{Verdict: REPLICATED.}
\end{abstract}

\maketitle

\section{Model}
Mn$_3X$ chiral AFMs carry a chiral spin order with negligible net dipole but a
finite cluster \emph{octupole} moment. Following the paper, the six-spin kagome
unit cell is coarse-grained to three sublattice magnetizations
$\vec M_i=M_s\vec m_i$ with free-energy density
\begin{equation}
F=\sum_{i\neq j}\!\left[J_{ij}\,\vec m_i\!\cdot\!\vec m_j
+D_{ij}\,\hat z\!\cdot\!(\vec m_i\!\times\!\vec m_j)\right]
-\sum_i K_i(\vec m_i\!\cdot\!\vec e_i)^2 ,
\end{equation}
with $J\!>\!D\!\gg\!K$, and a strain symmetry-breaking $J_{12}=J(1-\nu)$,
$\nu\ll1$. The octupole unit vector is
$\vec m=M_{xz}[\vec m_3+R\vec m_1+R^2\vec m_2]/3$ (Eq.~2 of the paper), with
$R$ a $2\pi/3$ rotation. Integrating out the fast modes yields the effective
low-energy octupole Hamiltonian (paper Eq.~B19)
\begin{equation}
VF_{\mathrm{oct}}=\tfrac{3}{2}M_sH_JV\,m_z^2+\Delta\sin^2\phi_{\mathrm{oct}},
\label{eq:Foct}
\end{equation}
where $(m_z,\phi_{\mathrm{oct}})$ are canonically conjugate,
$H_J=(9J+3\sqrt3 D)/3M_s\sim100$~T is the exchange field playing the role of
the $\sim1$~T dipole field of an XY ferromagnet,
$\Delta=2\nu KJV/(J+\sqrt3 D)$ is the in-plane barrier, and
$H_K=2\Delta/3M_sV$.

\paragraph{Octupole operator provenance.}
The rank-3 cluster-octupole moment is represented with the symmetrized Stevens
operator $T_{xyz}$ built by the shared kernel
\texttt{ollie\_multipolar\_stevens\_landau\_kernel.py}; we verify it is
Hermitian, traceless, with trace-norm $0.866$ (provenance credit to the shared
kernels cache).

\section{Method}
\textbf{(1) Low-barrier ($\Delta\!\ll\!kT$).} Here octupole rotations feel no
barrier; correlations decay by precessional dephasing. Thermal fields excite
$m_z$ with Boltzmann weight from Eq.~\eqref{eq:Foct}, and each precesses at
$\gamma H_J m_z$, giving (paper Eq.~9)
\begin{equation}
C(t)=\big\langle\cos(\gamma H_J m_z t)\big\rangle_{\mathrm{Boltz}}
=\exp\!\left(-\tfrac12\omega_J^2t^2\right),
\end{equation}
a Gaussian decay. We evaluate $C(t)$ by direct Monte-Carlo sampling of $m_z$
and extract $\tau$ at $C=1/e$. The self-consistent closed form for our free
energy is $\tau_{\mathrm{model}}=\sqrt2/\omega_J$,
$\omega_J=\gamma\sqrt{H_JH_{th}/3}$, $H_{th}=kT/M_sV$; the paper's published
prefactor (Eq.~10) is $\tau=\sqrt{2\ln2}/(\gamma\sqrt{H_JH_{th}})$. Both differ
only by a fixed dimensionless factor $\sim2.08$ and yield the identical ps
scale.

\textbf{(2) High-barrier ($\Delta\!>\!kT$).} Relaxation is thermal escape over
the barrier (random-telegraph noise). We use Langer's intermediate-to-high
damping formula (paper Eq.~D12)
\begin{equation}
\tau_{\mathrm{ihd}}=\frac{4\pi}{\gamma H_J\!\left[\sqrt{(\alpha(1+h_p))^2+4h_p}
-\alpha(1-h_p)\right]}\,e^{\Delta/kT},
\end{equation}
$h_p=H_K/H_J$, with octupole $\tau=0.25\,\tau_{\mathrm{esc}}$ and depopulation
factor $A\!\approx\!1$ for Mn$_3$Sn.

\textbf{(3) Direct Langevin.} We integrate the reduced stochastic dynamics
$\dot\phi=\gamma H_Jm_z$,
$\dot m_z=-\gamma H_K\sin\phi\cos\phi-\alpha\gamma H_Jm_z+\xi(t)$ with FDT noise
($\mathrm{var}\,m_z=kT/3M_sH_JV$), and measure
$C(t)=\langle m_x(0)m_x(t)\rangle$, $m_x=\cos\phi_{\mathrm{oct}}$.

\section{Results}
\begin{table}[t]
\caption{Low-barrier octupole relaxation time vs.\ dot volume ($T=300$~K,
$H_J=100$~T). Model = self-consistent 1/e time; paper-Eq10 = published
prefactor; numeric = Monte-Carlo of Eq.~9.}
\begin{ruledtabular}
\begin{tabular}{cccc}
$V$ (nm$^3$) & $\tau_{\mathrm{model}}$ (ps) & $\tau_{\mathrm{Eq10}}$ (ps)
& $\tau_{\mathrm{numeric}}$ (ps)\\\colrule
300    & 13.5 & 6.5  & 13.5\\
1000   & 24.6 & 11.8 & 24.7\\
3000   & 42.7 & 20.5 & 42.9\\
10000  & 77.9 & 37.5 & 77.8\\
30000  & 135.0& 64.9 & 135.1\\
100000 & 246.4& 118.5& 246.6\\
\end{tabular}
\end{ruledtabular}
\end{table}

The numeric Monte-Carlo evaluation of the dephasing integral matches the
self-consistent analytic form to a mean relative difference of $0.16\%$. The
smallest-volume (sub-$kT$) point reaches $\tau_{\mathrm{model}}=13.5$~ps
($\tau_{\mathrm{Eq10}}=6.5$~ps)---squarely the paper's ``$\sim$10~ps''
claim. The independent Langevin simulation at $V=10^3$~nm$^3$,
$\Delta=0.05\,kT$ gives $\tau=37.5$~ps vs.\ model $24.6$~ps (ratio $1.52$),
independently confirming the ps--tens-of-ps scale. In the high-barrier regime
($\Delta=4.5\,kT$, $\alpha=10^{-3}$) we obtain $\tau\approx13$~ns---three-plus
orders of magnitude slower, reproducing the two-mechanism crossover.

\section{Comparison and Verdict}
\begin{itemize}
\item ps-scale reached: \textbf{yes} (min $\tau\approx6$--$13$~ps).
\item Eq.~9 Monte-Carlo vs.\ analytic model: 0.16\% mean rel.\ diff.
\item Direct Langevin vs.\ analytic: within $1.5\times$.
\item High- vs.\ low-barrier ordering: escape (ns) $\gg$ dephasing (ps).
\end{itemize}
All four criteria pass. \textbf{Verdict: REPLICATED} (Coverage 9/10,
Agreement 8/10). The one honest caveat is a fixed $\sim$2$\times$
dimensionless prefactor between our self-consistent 1/e time and the paper's
published Eq.~10 constant, stemming from the mz-mode normalization and the
1/e-vs-1/2 crossing convention; it does not affect the physical scale or the
headline claim.

\section*{Provenance}
Octupole operators via
\texttt{ollie\_multipolar\_stevens\_landau\_kernel.py} (shared-kernels-cache).
All numbers computed at run time by
\texttt{konakanchi2025\_replication.py}; no values fabricated.

\end{document}
